% Source-derived chapter 04: Inputs from crystalline and prismatic Dieudonn\'{e} theory
% Original source: purity-flat-cohomology
\chapter{Inputs from crystalline and prismatic Dieudonn\'{e} theory}
\label{purity-flat-cohomology-chapter-04}
\source{purity-flat-cohomology}

\begin{remark}[Source section]
\label{purity-flat-cohomology-chapter-04-source}
\source{purity-flat-cohomology}
\source{purity-flat-cohomology}
This chapter reproduces the corresponding section of the retrieved
LaTeX source, including its definitions, statements, examples, and
proofs. It has not yet been translated into Lean.
\notready
\end{remark}

% The source section title was: Inputs from crystalline and prismatic Dieudonn\'{e} theory
\label{chapter:Dieudonne-theory}
\source{purity-flat-cohomology}


Our eventual source of the characteristic-primary aspects of purity for flat cohomology is a relation to coherent cohomology and the vanishing of the latter in presence of enough depth. To exhibit this relation, we use crystalline and prismatic Dieudonn\'{e} theories that classify commutative, finite, locally free groups of $p$-power order over perfect and perfectoid rings in terms of Dieudonn\'{e} modules. We review the crystalline classification in \S\ref{section:Dieudonne-positive-characteristic} and its prismatic generalization in \S\ref{section:Dieudonne-mixed-characteristic}. %, and, crucially for Chapter \ref{reduction}, establish the positive characteristic case of the key formula \eqref{eqn:key-formula} in \Cref{KT-input}.


%The former already gives the positive characteristic 

%the promised cohomological relation in positive characteristic, see . This case will be , which, in turn, will allow us to then deduce such a relation in mixed characteristic in \Cref{thm:perfectoidsupport}.






%%%%%%%%%%%%%%%%%%%%%%%%%%%%%%%%

\csub[Finite, locally free groups of $p$-power order over perfect rings] \label{section:Dieudonne-positive-characteristic}
\source{purity-flat-cohomology}

The positive characteristic case of the key formula \eqref{eqn:key-formula} is a perfect ring variant of Kato--Trihan's \cite{KT03}*{Proposition~5.10}. To establish it in \Cref{KT-input} we build on  Gabber's suggestion to use the pro-fppf site (see \Cref{thm:KTG-input}; alternatively, one could adapt the arguments of \emph{op.~cit.}). A key input %to both approaches 
is the crystalline classification of commutative, finite, locally free groups of $p$-power order over perfect $\bF_p$-schemes due to Berthelot, Gabber, and Lau, which we now review.

%, especially, \eqref{crystalline-cor}, which fairly directly implies most of the remaining positive characteristic cases of \Cref{main} (see \Cref{thm:main-pos-char} for a precise statement) and which will be a crucial input to Chapter \ref{reduction}. 
%In equal characteristic, relations between flat cohomology and coherent cohomology that rest on crystalline Dieudonn\'{e} theory have already been used in \cite{KT03}. Our approach builds on the one used there: more precisely, our 
%Our \Cref{KT-input} is a perfect scheme version of  The proof  and rests on the  

%\me{Put in the formula in terms of $\Ext^1$}

\bpp[Crystalline Dieudonn\'{e} modules over perfect bases] \label{Gab-eq}
\source{purity-flat-cohomology}
For a perfect $\bF_p$-scheme $S$, by an unpublished result of Gabber that built on \cite{Ber80}*{th\'{e}or\`{e}me 3.4.1} and was reproved by Lau \cite{Lau13}*{Corollary 6.5} by a different method, there is a covariant equivalence of categories
\[
G \mapsto \bM(G)
\]
from the category of commutative, finite, locally free $S$-groups $G$ that are locally on $S$ of $p$-power order to the category of quasi-coherent %, (locally on $S$) finitely presented 
$W(\sO_S)$-modules $\bM$ equipped with Frobenius (resp.,~inverse Frobenius) semilinear maps $F\colon \bM \ra \bM$ (resp.,~$V\colon \bM \ra \bM$) with $FV = VF = p$ such that for every  affine open $\Spec(R) \subset S$ the $W(R)$-module $\Gamma(R, \bM)$ is finitely presented, killed by a power of $p$, and of projective dimension $\le 1$. The functor is defined by Zariski-local glueing as follows (see \cite{Lau13}*{proof of Corollary 6.5}): Zariski locally on $S$ one finds $p$-divisible groups $H_0$, $H_1$ that fit into an exact sequence
\[
0 \ra G \ra H_0 \ra H_1 \ra 0 \qq \x{and sets} \qq \bM(G) \ce \Coker(\bM(H_0) \hra \bM(H_1)),
\]
where $\bM(H_i) \ce \Gamma((S/\bZ_p)_\cris, \bD(H_i))$ is the evaluation of the covariant Dieudonn\'{e} crystal 
\[
\bD(H_i) \ce \sE xt^1_{(S/\bZ_p)_\cris}(H_i^*, \sO_{(S/\bZ_p)_\cris}), 
\]
which, by \cite{BBM82}*{section 5.3}, is identified with,
 \[
(\sE xt^1_{(S/\bZ_p)_\cris}(H_i, \sO_{(S/\bZ_p)_\cris}))^* 
\]
(the dual of the locally free crystal of $H_i$ defined in \cite{BBM82}*{d\'{e}finition~3.3.6, th\'{e}or\`{e}me~3.3.10}) at the terminal ind-object $\{(S, W_n(\sO_S))\}_{n \ge 0}$ of the crystalline site $(S/\bZ_p)_\cris$. %; the resulting $\bM(G)$ does not depend on the choice of $H_1 \ra H_2$, and one globalizes it by glueing. 
Since \emph{op.~cit.}~uses big crystalline sites, the formation of $\bM(G)$ commutes with base change to any perfect $S$-scheme. 
\epp



\beg \label{Gab-eq-ex}
\source{purity-flat-cohomology}
By \cite{BBM82}*{exemple~4.2.16~(i)}, the $W(\sO_S)$-module that underlies $\bb1_n \ce \bM(\bZ/p^n\bZ)$ is
\[
W(\sO_S)/p^n \qxq{with} F = p\cdot\Frob(-), \q V = \Frob\i(-).
\]
Likewise, the $W(\sO_S)$-module that underlies $\bM(\mu_{p^n})$ is  
\[
W(\sO_S)/p^n \qxq{with} F = \Frob(-), \q V = p\cdot\Frob\i(-).
\] 
\eeg



\blem \label{MG-exact}
\source{purity-flat-cohomology}
In \uS\uref{Gab-eq}, the functor $G \mapsto \bM(G)$ and its inverse preserve short exact sequences. 
\elem

\bpf
By \cite{BBM82}*{proposition~1.1.7, th\'{e}or\`{e}me~3.3.3}, the functor $H_i \mapsto \bD(H_i)$ preserves short exact sequences, hence so does $H_i \mapsto \bM(H_i)$ (compare with \cite{BBM82}*{proposition~1.1.19, d\'{e}finition~1.2.1}). % Left-exactness is clear and if the surjection wouldn't persist, then the cokernel would propagate by base change (tensor product is right-exact), which would contradict the exactness of $H_i \mapsto \bD(H_i)$. I think those ``compare with'' references are just for moral support. Also, the following is a much worse argument that I had initially.
%Moreover, the forgetful morphism $u \colon (S/\bZ_p)_\cris \ra S_\et$ of topoi satisfies $R^iu_*(D) = 0$ for every $i > 0$ and every crystal $D$, as follows from the perfectness of $S$ and the identification of $Ru_*(D)$ with the de Rham complex of a module with a connection determined by $D$ (for instance, one may apply the much more general \cite{Bei13a}*{1.8.1}). 
To deduce the same for the functor $G \mapsto \bM(G)$, thanks to the snake lemma, it suffices to Zariski locally on $S$ embed a short exact sequence $0 \ra G \ra G' \xra{\pi} G'' \ra 0$ into a short exact sequence $0 \ra H_0 \ra H_0' \ra H_0'' \ra 0$ of $p$-divisible groups (see \cite{BBM82}*{lemme~3.3.12}). For this, we choose Zariski local embeddings into $p$-divisible groups $\iota'\colon G' \hra H_0$ and $\iota'' \colon G'' \hra H_0''$ and replace $\iota'$ by 
\[
(\iota',\, \iota''\circ \pi)\colon G' \hra  H_0 \times H_0'' \equalscolon H_0'.
\] 
%to arrange a compatible with $\pi$ surjection $H_0' \surjects H_0''$
%finding $H_0' \surjects H_0''$ suffices. 
%, consider the 
% into a $p$-divisible group $H$ and apply the snake lemma to the resulting short exact sequence morphism 
%\[
%\xymatrix{
%0 \ar[r] & \bM(H) \ar@{^(->}[d]\ar[r] & \bM(H/G') \ar@{=}[d]\ar[r] & \bM(G') \ar[r] \ar@{->>}[d] & 0 \\
%0 \ar[r] & \bM(H/G) \ar[r] & \bM(H/G') \ar[r] & \bM(G'') \ar[r] & 0
%}
%\]
For the remaining exactness of the inverse, granted that 
\[
0 \ra \bM(G_1) \ra \bM(G_2) \ra \bM(G_3) \ra 0
\]
is a short exact sequence, we need to show that the complex $G_1 \ra G_2 \ra G_3$ is also a short exact sequence. In the case when $S$ is geometric point we may decompose this complex into short exact sequences of finite flat group schemes and conclude by the exactness of $G \mapsto \bM(G)$. Thus, in general we check on   geometric $S$-fibers that $G_1 \isomto \Ker(G_2 \ra G_3)$ (see \cite{EGAIV4}*{corollaire 17.9.5}). It then remains to note that $G_2/G_1 \hra G_3$ becomes an isomorphism after applying $\bM(-)$, so is an isomorphism.
\epf




The equivalence $G \mapsto \bM(G)$ leads to the following description of the low degree cohomology of $G$. 

\bprop \label{lem:syntomic-low-coh-id}
\source{purity-flat-cohomology}
Let $S$ be a perfect $\bF_p$-scheme and let $G$ be a commutative, finite, locally free $S$-group killed by $p^n$. 
%In \uS\uref{synt-cx}, if the order of $G$ is constant on $S$ and equal to $p^n$, then 
 We have the following functorial in $G$ identifications of sheaves on $S_\et$\ucolon
\[\ba \label{H0-cris-id} %\label{H1-cris-id}
\source{purity-flat-cohomology}
G\cong \,\sH om_S(\bZ/p^n\bZ, G) \overset{\ref{Gab-eq}}{\cong}  \sH om_{W_n(\sO_S),\,F,\, V}(\bb1_n, \bM(G)) \cong \bM(G)^{V = 1}, \\
\sE xt^1_S(\bZ/p^n\bZ, G) \overset{\ref{Gab-eq},\, \ref{MG-exact}}{\cong}  \sE xt^1_{W_n(\sO_S),\, F,\, V}(\bb1_n, \bM(G)) \cong \bM(G)/(V - 1)(\bM(G)),
\ea \]
where $\sE xt^1_S$ denotes the \'{e}tale sheafification of the functor of extensions of fppf $\bZ/p^n\bZ$-module sheaves.
\eprop

\bpf
The full faithfulness of $G \mapsto \bM(G)$ gives the first line of the displayed identification: the map to $\bM(G)^{V = 1}$ is the evaluation $(f \colon \bb1_n\, \ra\, \bM(G))\, \mapsto\, f(1)$. For the second line, we define the last identification as follows. To a local section $m$ of $\bM(G)$, we associate the extension $\bM(G) \oplus \bb1_n$ for which the Verschiebung is determined by $(0, 1) \mapsto (m, 1)$ and the Frobenius is then necessarily determined by $(0, 1) \mapsto (-F(m), p)$ (we write $F$ and $V$ for those of $\bM(G)$). Such extensions for $m$ and $m'$ are isomorphic if and only if the isomorphism of $W_n(\sO_S)$-module extensions determined by $(0, 1) \mapsto (a, 1)$ for some $a \in \bM(G)$ is $V$- and $F$-equivariant. The $V$-equivariance amounts to $(m + a, 1) = (m' + V(a), 1)$, that is, to $m - m' \in (V - 1)(\bM(G))$, and the $F$-equivariance amounts to $(-F(m) + pa, p) = (-F(m') + F(a), p)$ and follows from $V$-equivariance. Since any extension of $\bb1_n$ by $\bM(G)$ is \'{e}tale locally split as an extension of $W_n(\sO_S)$-modules, the claim follows.
\epf




In \Cref{KT-input}, we will upgrade the identifications of \Cref{lem:syntomic-low-coh-id} to a formula that expresses the flat cohomology $R\Gamma(S, G)$ in terms of the quasi-coherent cohomology $R\Gamma(W_n(S), \bM(G))$. For this, we first show in \Cref{thm:KTG-input} that the map $V - 1 \colon \bM(G) \ra \bM(G)$ is pro-fppf locally surjective.

%The proof of \Cref{KT-input} uses the pro-fppf local description of $G$ presented .


\bpp[The pro-fppf site] \label{X-profppf}
\source{purity-flat-cohomology}
A scheme map $X' \ra X$ is \emph{pro-ppf} (resp.,~\emph{pro-fppf}) if $X'$ may be covered by opens $\Spec(A') \subset X'$ for which there is a factorization 
\[
\Spec(A') \ra \Spec(A) \subset X
\]
in which the $A$-algebra $A'$ is a filtered direct limit of flat, finitely presented $A$-algebras (compare with \S\ref{ind-fppf-maps}) and $\Spec(A) \subset X$ is an open immersion (resp.,~and $X' \ra X$ is faithfully flat). Pro-ppf maps 
%\[%be \label{profppf-cover}
$\{X_i' \ra X\}_{i \in I}$ 
%\ee
form a \emph{pro-fppf cover} if each quasi-compact open of $X$ is a finite union of images of quasi-compact opens of $\bigsqcup_{i \in I} X_i'$.  By \S\ref{ind-fppf-maps}, the pro-ppf maps are stable under base change and composition, so the category of $X$-schemes with pro-fppf covers as coverings defines the \emph{pro-fppf site} of $X$. A pro-fppf cover is also an fpqc cover, so an fpqc sheaf is also a pro-fppf sheaf.
\epp



\blem \label{lem:perf-still-pro-fppf}
\source{purity-flat-cohomology}
Let $S$ be a perfect $\bF_p$-scheme and let $X^\perf$ denote the perfection of an $\bF_p$-scheme $X$. If $\{X_i \ra S\}_{i \in I}$ is an fpqc \up{resp.,~pro-fppf} cover of $S$, then so is $\{X_i^\perf \ra S\}_{i \in I}$.
%For a pro-fppf cover  of a perfect $\bF_p$-scheme $S$, the termwise perfection  is also a pro-fppf cover of $S$.
%of a perfect  can be defined by a pro-fppf cover $\{X'_i \ra S\}$
%pro-fppf site of a perfect $\bF_p$-scheme has a base consisting of perfect schemes. 
\elem

\bpf
The composite 
\[
X_i \xra{\Frob^n} X_i \ra S \qxq{factors as} \xymatrix@C=10pt{X_i \ar[r] &S \ar[rr]^-{\Frob^n}_-{\sim} &&S},
\]
so it is fpqc (resp.,~pro-fppf). Thus, the inverse limit $X_i^\perf$ of 
\[
\dotsc \xra{\Frob} X_i \xra{\Frob} X_i
\]
is fpqc (resp.,~pro-fppf) over $S$ (see \S\ref{ind-fppf-maps}).
%with affine transition maps is an fpqc (resp.,~pro-fppf) cover of $S$. 
\epf





\bprop \label{thm:KTG-input}
\source{purity-flat-cohomology}
Let $S$ be a perfect $\bF_p$-scheme and let $G$ be a commutative, finite, locally free $S$-group that is locally on $S$ of $p$-power order. There is a functorial in $G$ short exact sequence 
\be \label{eqn:KTG-input-SES}
\source{purity-flat-cohomology}
0 \ra G \xra{\ref{lem:syntomic-low-coh-id}} \bM(G) \xra{V - 1} \bM(G) \ra 0
\ee
of sheaves on the category of perfect $S$-schemes endowed with the pro-fppf topology. %,  where the map $G \ra \bM(G)$ is given by applying $\bM(-)$ to $\Hom_S(\underline{\bZ_p}_S, G)$ \up{which locally on $S$ is $\Hom_S(\underline{\bZ/p^n\bZ}_S, G)$}.
\eprop

\bpf
The left exactness follows from \Cref{lem:syntomic-low-coh-id}. For the remaining surjectivity we may work \'{e}tale locally on $S$, so, by \Cref{lem:syntomic-low-coh-id}, we need to show that a given extension 
\[
0 \ra G \ra G' \xra{\pi} \underline{\bZ/p^n\bZ}_S \ra 0
\]
of $\bZ/p^n\bZ$-module sheaves splits over a pro-fppf cover of $S$. The extension splits over the fppf cover $\pi\i(1) \surjects S$, which is a $G$-torsor, %_{\underline{\bZ/p^n\bZ}_S}$-torsor), 
so \Cref{lem:perf-still-pro-fppf} supplies a desired pro-fppf cover. (The point of using \Cref{lem:perf-still-pro-fppf}, so, relatedly, the pro-fppf rather than simply the fppf topology, is to stay in the realm of perfect base schemes in order to be able to consider the functor $\bM(-)$.)
\epf


%We are ready for the promised relation between flat and quasi-coherent cohomologies.

\bthm \label{KT-input}
\source{purity-flat-cohomology}
Let $S$ be a perfect $\bF_p$-scheme and let $G$ be a commutative, finite, locally free $S$-group that locally on $S$ is of $p$-power order. The map  of sites $\eps \colon S_\fppf \ra S_{\et}$ gives a functorial triangle
\be \label{eqn:KT-triangle}
\source{purity-flat-cohomology}
R\eps_*(G) \ra \bM(G) \xra{V - 1} \bM(G) \ra (R\eps_*(G))[1] \qxq{on} S_\et.
\ee
%\up{in other words, $R\eps_*(G) \simeq \sS_G$, see \uS\uref{synt-cx}}\uscolon 
In particular, if $G$ is killed by $p^n$, then, %, in addition, $S = \Spec(A)$ is affine \up{with $A$ perfect}, then 
for any closed subset $Z \subset S$, we have
\be \label{crystalline-cor}
\source{purity-flat-cohomology}
R\Gamma_Z(S, G) \cong R\Gamma_Z(W_n(S), \bM(G))^{V = 1}  %\qxq{and} H^i(A, G) = 0 \qxq{for} i \ge 2.
\ee
functorially in $S$, $Z$, and $G$.
\ethm




%The following argument for \Cref{KT-input} is based on a suggestion of Gabber.


\bpf
The identification \eqref{crystalline-cor} follows from \eqref{eqn:KT-triangle} by applying $R\Gamma_Z(S, -)$. For the latter, we fix a suitable auxiliary %strong limit cardinal of uncountable cofinality 
cutoff cardinal $\kappa$ with 
$\kappa > \abs{S}$ (see \S\ref{conv}), consider the small pro-fppf site $S_{\profppf,\, \kappa}$ bounded by $\kappa$, its subsite $S_{\profppf,\, \perf,\, \kappa}$ of perfect schemes, and the morphisms
\[
\xymatrix@C=15pt@R=12pt{
S_{\profppf,\, \kappa} \ar[r]_-b \ar[d]_-a & S_{\profppf,\, \perf,\, \kappa} \ar[d]_-c \\
S_\fppf \ar[r]^-{\eps} & S_\et.
}
\]
By limit arguments, %the functor $S' \mapsto R\Gamma_\fppf(S', G)$ satisfies pro-fppf descent, so, since its values are bounded below, it also satisfies pro-fppf hyperdescent. %(see, for instance, \cite{CM19}*{2.5}). 
%It follows that 
$R^{\ge 1}a_*(G) \cong 0$, so 
\[
R\eps_*(G) \cong R(\eps \circ a)_*(G) \cong R(c \circ b)_*(G) \cong Rc_*(Rb_*(G)).
\]
By \Cref{lem:perf-still-pro-fppf}, the functor $b_*$ is exact, so $R\eps_*(G) \cong Rc_*(G)$. Moreover, by %Similarly to the argument for $a_*$, 
faithfully flat descent for modules (see \cite{SP}*{Lemma~\href{https://stacks.math.columbia.edu/tag/023M}{023M}}), we have 
\[
R^{\ge 1}c_*(\bM(G)) \cong 0.
\]
% To see this "by faithfully flat descent" one argues as follows. Firstly, by applying $W_n(-)$ to a flat map of perfect rings one gets a flat map due to the infinitesimal criterion of flatness. Secondly, $W_n(-)$ commutes with tensor products of perfect rings, see (2.1.10.1). Granted this, applying $W_n(-)$ is harmless and one is really studying quasi-coherent cohomology. Now to show that $R^nc_*(\bM(G)) = 0$ we argue by induction on $n$ (for all $S$ simultaneously). For $n = 1$ it follows by flat descent by taking finer and finer pro-fppf covers. Granted that the claim is known for $i < n$ (and all $S$) one can use the Leray spectral sequence to reduce to checking exactness at the $n$-th spot of the complex as in the SP reference (we apply the inductive assumption to self-products of the cover to get vanishing of $H^{< n}$ there). Due to this argument, there is no need to do hypercovers! (Although we also have a reference for hypercover flat descent for modules, in some appendix of SAG, I think it is referenced in Chapter 5). 
Thus, by applying $Rc_*(-)$ to the short exact sequence \eqref{eqn:KTG-input-SES} we get the desired \eqref{eqn:KT-triangle} (independently of the choice of $\kappa$). % If you enlarge $\kappa$, then you get an extra row on top of the commutative diagram and by chasing through see that $Rc_*(G)$ remains the same.  
\epf








%\kestutis{I wrote it this way in order to make the triangle functorial in $S$. }


%\kestutis{
%There are difficulties in the proof caused by the nonrepleteness of the involved topoi. Thus, should instead formulate this for the pushforward along $S_{\rm{profppf}} \ra S_{\rm{proet}}$. Here are some technical points:
%\benuma
%\item
%The profppf site is defined by imitating \cite{Sch17}*{7.8 and 8.1 (i)} and using an uncountable strong limit cardinal $\kappa$ to bound the sizes in order to avoid set theoretic issues. It gives a replete topos: argue as for the fpqc topos in \cite{BS15}*{3.1.7} using the criterion \cite{SP}*{\href{https://stacks.math.columbia.edu/tag/07C3}{07C3}} to recognize ind-fppf algebras (e.g., $\kappa$-filtered direct limits of affines are again objects of the site). There are two version of the site: the big and small one; officially, we use the big one to get better topos functoriality, but in practice and for arguments we use the small one; the cohomology is the same by using \cite{Mil80}*{III.3.1} (SGA4 reference?).

%\item
%For all sheaves $\sF$ on $S_\fppf$ and the map of sites $\nu \colon S_{\rm{profppf}} \ra S_\fppf$ we have $\sF \isomto R\nu_*(\nu^*(\sF))$. This is proved the same way as in \cite{BS15}*{5.1.6} using limit arguments and is quite formal (key input: limit formalism holds for fppf cohomology).

%\item
%The pro\'{e}tale site $S_\proet$ is defined not as in \cite{BS15} but as in \cite{Sch17}*{7.8 and 8.1 (i)}. This is needed in order to make sense of $S_{\rm{profppf}} \ra S_{\rm{proet}}$. The topos does not change by \cite{SGA4I}*{III.4.1} and \cite{BS15}*{2.3.4}. It is replete by \cite{BS15}*{4.2.8, 3.2.3}, which show that it is even locally contractible (or argue directly as in (1)).

%\item
%The formation of $\bD(G)$ commutes with arbitrary base change because it is defined using the big crystalline site. Then also the formation of $\bM(G)$ commutes with base change ($p$-completed in the Witt vector direction). By a limit argument, ind-\'{e}tale morphisms have relative Frobenius to be an isomorphism. Thus, perfectness of an $\bF_p$-scheme is an ind-\'{e}tale local condition. Thus, perfectness ascend along weakly \'{e}tale morphisms due to \cite{BS15}*{2.3.4}. (The items in this point are only tangentially relevant though.)

%\item
%Get the analogue of \Cref{Gro-lem} by reducing to the fppf vs \'{e}tale statement (i.e., to \Cref{Gro-lem}).
%\eenum
%}



%\brem 
%Since $G \tensor^\bL_{\bZ} \bZ/p^n\bZ \cong (G[p^n])[1]$, we have $G[p^N] \tensor^\bL_{\bZ/p^N\bZ} \bZ/p^n\bZ \cong G[p^n]$ whenever $N \ge n$, so that, by the projection formula \cite{SP}*{\href{https://stacks.math.columbia.edu/tag/0944}{0944}}, also $R\eps_*(G[p^N]) \tensor^\bL_{\bZ/p^N\bZ} \bZ/p^n\bZ \cong R\eps_*(G[p^n])$. Thus, by passing to the limit in $N$, we get $(R\lim(R\eps_*(G[p^n]))) \tensor^\bL_{\bZ} \bZ/p^n\bZ \cong R\eps_*(G[p^n])$, to the effect that \Cref{KT-input} supplies the following exact triangle for every $n \ge 0$:
%\be\label{mod-p-triangle}
%R\eps_*(G[p^n]) \ra \bM(G)^0/p^n \xra{1 - F/p} \bM(G)/p^n \ra R\eps_*(G[p^n])[1].
%\ee
%\erem

%
%The alternative proof of \Cref{KT-input} given in \S\ref{pp:KT-fix} uses the following well-known lemma.
%

%
%\bpp[Another proof for \Cref{KT-input}] \label{pp:KT-fix}
%The following alternative argument for \Cref{KT-input} adapts that for \cite{KT03}*{5.10}. We include it because %we believe that 
%it may be useful for fully understanding \cite{KT03}*{5.10} (the direct limits in its argument do not appear to be filtered; see also \cite{TV18}*{1.0.1} for a discussion of other difficulties there that originate in the nonrepleteness of fppf topoi). % Secretly: I do not fully understand the argument given there (in their setting, the definition of $\sS_D$ is a little bit troubling: seems like lack of functoriality of the mapping cone could be a problem in really getting a canonical identification). Also, their system of $G' \ra G''$ is not filtered I think. According to arXiv preprints of Trihan and Vauclair it seems that there unfortunately are serious difficulties with \cite{KT03}*{5.10}.
%
%For a perfect $\bF_p$-scheme $S$ and a commutative, finite, locally free $S$-group $G$ locally on $S$ of $p$-power order, the \emph{syntomic complex} of $G$ is  the concentrated in degrees $0$ and $1$ complex 
%\be \label{synt-def}
%\sS_G \ce [\bM(G) \xra{V - 1} \bM(G)] \qx{of sheaves on $S_\et$},
%\ee
%so that $\sS_G$ is also identified with the mapping fiber (that is, with $\Cone(*)[-1]$) of $V - 1$ considered as a map between sheaves placed in degree $0$. %For example, $\sS_{\bb1} \cong [W(\sO_S) \xra{1 - \Frob(-)} W(\sO_S)]$. 
%By \S\ref{Gab-eq}, a short exact sequence $0 \ra G \ra G' \ra G'' \ra 0$ of commutative, finite, locally free $S$-groups of $p$-power order induces a short exact sequence 
%\be \label{synt-ex}
%0 \ra \sS_G \ra \sS_{G'} \ra \sS_{G''} \ra 0.
%\ee
%The formation of the complex $\sS_G$ commutes with base change to any perfect $S$-scheme. 
%
%
%In \Cref{KT-input}, our goal is to show that $R\eps_*(G) \simeq \sS_G$. By \eqref{H0-cris-id} and \Cref{Gro-lem}, the cohomologies of $R\eps_*(G)$ and $\sS_G$ agree, and the idea is to express $R\eps_*(G)$ in terms of homomorphisms between finite, locally free group schemes and then transfer to an analogous expression for $\sS_G$. For simplicity, in this proof we neglect the functoriality in $S$: at least among perfect $S$-schemes of bounded size, such functoriality may be obtained by carrying out the constructions below with affine perfect $S$-schemes of size $< \kappa$ in place of $S_{\et,\,\aff}$ for a sufficiently large strong limit cardinal $\kappa$. 
%
%By passing to a clopen subscheme of $S$, we may assume that the order of $G$ is globally constant and equal to $p^n$. Moreover, it suffices to exhibit the desired triangle after restricting to the full subsite $S_{\et,\, \aff}$ of $S_{\et}$ of affine schemes because this does not change the associated topos. % Then take the sheafification to get the triangle.
%By replacing $S_{\et,\, \aff}$ by an equivalent subsite, we assume that its collection of objects forms a set. The purpose of restricting to $S_{\et,\, \aff}$ is to be able to commute the fppf cohomology of objects in $S_{\et,\, \aff}$ with filtered direct limits of sheaves (for which one needs quasi-compactness and quasi-separatedness).
%
%For every object $S'$ of $S_{\et,\, \aff}$, we define a contravariantly functorial in $S'$ set $J_{S',\, \infty}$ of extensions 
%\[
%0 \ra G_{S'} \ra G' \ra G'' \ra 0
%\]
%of commutative, finite, locally free $S'$-group schemes killed by $p^n$ as follows. For each $S'$, we choose a set $J_{S',\, 0}$ of representatives for the set of isomorphism classes of such extensions. Then we define $J_{S',\, n}$ for $n > 0$ inductively by letting it be the disjoint union of the base changes of the elements of $J_{S'',\, n - 1}$ along variable morphisms $S' \ra S''$ in $S_{\et,\, \aff}$ with the fixed source $S'$. The identity morphism allows us to view $J_{S',\, n - 1}$ as a subset of $J_{S',\, n}$, and we obtain the promised $J_{S',\, \infty}$ by setting
%\[
%\tst J_{S',\, \infty} \ce \bigcup_{n \ge 0} J_{S',\, n}.
%\]
%We let $J_{S'}$ be the (functorial in $S'$) set of all finite subsets of $J_{S',\, \infty}$ and, for an element $\Sigma \in J_{S'}$, let
%\[
%0 \ra G_{S'} \ra G'_\Sigma \ra G''_{\Sigma} \ra 0
%\]
%be the pushout sum of the extensions in $\Sigma$. % This dry part is a little bit similar to \cite{Schr17}*{\S6}.
%Forming unions of subsets of $J_{S',\,\infty}$ filters the set $J_{S'}$, so we obtain a contravariantly functorial in $S'$ short exact sequence of sheaves on $S'_\fppf$:
%\[
%\tst 0 \ra G_{S'} \ra \dirlim_{\Sigma \in J_{S'}} (G'_\Sigma) \ra \dirlim_{\Sigma \in J_{S'}} (G''_{\Sigma}) \ra 0.
%\]
%Consequently, we obtain likewise functorial exact triangles (see \cite{SP}*{\href{https://stacks.math.columbia.edu/tag/0739}{0739}})
%\[
%\tst R\Gamma(S'_\fppf, G) \ra \dirlim_{\Sigma \in J_{S'}} R\Gamma(S'_\fppf, G'_\Sigma) \ra \dirlim_{\Sigma \in J_{S'}} R\Gamma(S'_\fppf, G''_{\Sigma}) \ra R\Gamma(S'_\fppf, G)[1].
%\]
%We claim that for every object $S'$ of $S_{\et,\, \aff}$ the following vertical morphisms
%\be\ba \label{magic-square}
%\xymatrix{
%\tst \dirlim_{\Sigma \in J_{S'}} \Hom_{S'}(\bZ/p^n\bZ, G'_{\Sigma}) \ar[r] \ar[d] &\dirlim_{\Sigma \in J_{S'}} \Hom_{S'}(\bZ/p^n\bZ, G''_{\Sigma}) \ar[d] \\
%\tst\dirlim_{\Sigma \in J_{S'}} R\Gamma(S'_\fppf, G'_{\Sigma}) \ar[r] & \tst \dirlim_{\Sigma \in J_{S'}} R\Gamma(S'_\fppf, G''_{\Sigma})
%}
%\ea \ee
%from the nonderived global sections induce an isomorphism on the $\tau_{\le 1}$ truncations of the mapping fibers of the horizontal arrows. For this, it suffices to note that the bottom horizontal map in \eqref{magic-square} is injective on degree $1$ cohomology: indeed, each element of 
%\[
%H^1(S'_\fppf, G) \cong \Ext^1_{S'}(\bZ/p^n\bZ, G_{S'})
%\]
%corresponds to a short exact sequence $0 \ra G_{S'} \ra G' \ra \bZ/p^n\bZ \ra 0$ of commutative, finite, locally free $S'$-groups killed by $p^n$, so dies in $H^1(S'_\fppf, G')$ and hence, by the construction of $J_{S'}$, also in $\dirlim_{\Sigma \in J_{S'}} H^1(S'_\fppf, G'_{\Sigma})$. Since sheafification is an exact triangulated functor, % here one uses the formalism of the presheaf topos as described in 4.1 of Ainf-coho.
%we conclude that the sheafification of the mapping fiber of the top horizontal map of \eqref{magic-square} is identified with the sheafification of $\tau_{\le 1}(R\Gamma(S'_\fppf, G))$, that is, with $\tau_{\le 1}(R\eps_*(G)) \overset{\ref{Gro-lem}}{\cong} R\eps_*(G)$. We will conclude by showing that this sheafification of the  mapping fiber is also the mapping fiber of $\bM(G) \xra{V - 1} \bM(G)$. 
%
%By \eqref{H0-cris-id}, the top horizontal map of \eqref{magic-square} is identified with that of the square
%\be \ba \label{more-magic}
%\xymatrix{
%\tst \dirlim_{\Sigma \in J_{S'}} \Gamma(S'_\et, H^0(\sS_{G'_{\Sigma}})) \ar[r] \ar[d] &\dirlim_{\Sigma \in J_{S'}} \Gamma(S'_\et, H^0(\sS_{G''_{\Sigma}})) \ar[d] \\
%\tst\dirlim_{\Sigma \in J_{S'}} R\Gamma(S'_\et, \sS_{G'_{\Sigma}}) \ar[r] & \tst \dirlim_{\Sigma \in J_{S'}} R\Gamma(S'_\et, \sS_{G''_{\Sigma}})
%}
%\ea \ee
%whose vertical maps are induced by the natural transformation $\tau_{\le 0}(R\Gamma(S'_\et, -)) \ra R\Gamma(S'_\et, -)$. Due to \eqref{H0-cris-id} and the argument that gave the injectivity of the bottom horizontal map of \eqref{magic-square} on degree $1$ cohomology, every element of $H^1(\sS_G)$ \'{e}tale locally on $S'$ dies in $\dirlim_{\Sigma \in J_{S'}} H^1(\sS_{G'_{\Sigma}})$. Thus, due to the hypercohomology spectral sequence, every element of $H^1(S'_\et, \sS_G)$ \'{e}tale locally on $S'$ dies in $\dirlim_{\Sigma \in J_{S'}} H^1(S'_\et, \sS_{G'_{\Sigma}})$. Consequently, the vertical maps of the square \eqref{more-magic} induce a map between the mapping fibers of the horizontal maps that is an isomorphism on degree $1$ cohomology after sheafification and an isomorphism on degree $0$ cohomology even before the sheafification. Thus, the sheafification of the mapping fiber of the top horizontal map of \eqref{more-magic} is $\tau_{\le 1}(\sS_G) \cong \sS_G$ and, due to the analogous conclusion about \eqref{magic-square}, we obtain the desired $R\eps_*(G) \simeq \sS_G$. \QED
%\epp
%
%




\bcor \label{cor:perfect-vanish}
\source{purity-flat-cohomology}
For a perfect $\bF_p$-algebra $A$ and a commutative, finite, locally free $A$-group $G$ of $p$-power order, 
\[
H^i(A, G) \cong 0 \qxq{for} i \ge 2.
\]
\ecor

\bpf
Since affines have no higher quasi-coherent cohomology, \eqref{crystalline-cor} with $Z = S$ suffices. %and the vanishing of quasi-coherent cohomology of affine schemes.
%The exact triangle \eqref{eqn:KT-triangle} reduces us to proving that $H^{\ge 1}(A, \bM(G)) = 0$. However, the $p$-adic filtration of $\bM(G)$ is finite, exhaustive, and its subquotients are quasi-coherent, so this vanishing follows from the vanishing of quasi-coherent cohomology of affine schemes.
\epf




\Cref{KT-input} implies much of the positive characteristic case of purity for flat cohomology. We record this in \Cref{thm:main-pos-char} because the intervening auxiliary lemmas are also important in the general case. The following example illustrates why positive characteristic is significantly simpler.



\beg \label{eg:depth-example}
\source{purity-flat-cohomology}
For a Noetherian local ring $(R, \fm)$, by the cohomological characterization of depth, $H^i_\fm(R, \bG_a) \cong 0$ for $i < \depth(R)$ (compare with \cite{SGA2new}*{expos\'{e} III, proposition~3.3 (iv)}). Thus, if $R$ is also an $\bF_p$-algebra, then the Frobenius-kernel and the Artin--Schreier sequences give %the following special cases of \Cref{main}: 
\be \label{eqn:depth-vanishing}
\source{purity-flat-cohomology}
H^i_\fm(R, \gA_p) \cong 0 \qxq{and} H^i_\fm(R, \bZ/p\bZ) \cong 0 \q \x{for} \q i < \depth(R).
\ee
Complete intersection $R$ have $\depth(R) = \dim(R)$, so for them \eqref{eqn:depth-vanishing} gives cases of \Cref{main}. 
%This, however, does not fully capture the importance  for general $G$, which  %, which also plays an important role in the general case of \Cref{main}.
% for complete intersections $R$ that are  and $G = \bZ/p\bZ$ or $G = \gA_p$ follows. 
\eeg


The complete intersection assumption fully manifests itself in the following lemma.
%The following lemma  It is also a place where the complete intersection assumption of the latter prominently manifests itself. 

\blem \label{mod-out-by-more}
\source{purity-flat-cohomology}
Let $A$ be a ring, let $f_1, \dotsc, f_m \in A$ be a regular sequence, and let $G$ be a commutative $(A/(f_1, \dotsc, f_m))$-group scheme that is either smooth or finite locally free. For all $n_1, \dotsc, n_m > 0$ and $f_1^{1/n_1}, \dotsc, f_m^{1/n_m} \in A$ and all ideals $I \subset A/(f_1, \dotsc, f_m)$ containing some $(A/(f_1, \dotsc, f_m))$-regular sequence $a_1, \dotsc, a_d$ of length $d$, the map
\[
H^i_I(A/(f_1, \dotsc, f_m), G) \ra H^i_I(A/(f_1^{1/n_1}, \dotsc, f_m^{1/n_m}), G) 
\]  %the map
is injective for $i < d$ and bijective for $i < d - 1$.
\elem


\bpf
We use the B\'{e}gueri sequence \eqref{Begueri-0} and the five lemma to assume that $G$ is smooth.
By \cite{SP}*{Lemma~\href{https://stacks.math.columbia.edu/tag/07DV}{07DV}}, the sequence $f_1^{i_1/n_1}, \dotsc, f_m^{i_m/n_m}, a_1, \dotsc, a_d$ is regular for all $i_1, \dotsc, i_m \ge 1$. Thus, if $n_1 > 1$, then, by induction on $m$,  the square-zero ideal $J$ that cuts out the closed immersion 
\[
j \colon \Spec(A/(f_1^{(n_1 - 1)/n_1}, f_2, \dotsc, f_m)) \hra \Spec(A/(f_1, f_2, \dotsc, f_m))
\]
is free as an $(A/(f^{1/n_1}_1, f_2, \dotsc, f_m))$-module.
%we may assume that $m = 1$ and set $f \ce f_1$ with $n \ce n_1$.  
%For the latter, we use \'{e}tale cohomology and consider  whose ideal $J$ 
 Moreover, deformation theory supplies the short exact sequence of \'{e}tale sheaves% or \Cref{prop:defthy} below)
\[%be \label{MOBM-eq}
\source{purity-flat-cohomology}
0 \ra \sH om (e^*(\Omega^1_{G/(A/(f_1,\, \dotsc,\, f_m))}), J) \ra G \ra j_*(G_{A/(f^{(n_1 - 1)/n_1}_1,\, f_2,\, \dotsc,\, f_m)}) \ra 0,
\]%ee
where the $\sH om$ is of $A/(f_1,\, \dotsc,\, f_m )$-module sheaves and $e$ is the unit section of $G$ (see, for instance, \cite{topology-torsors}*{Lemma~B.14 and its proof}). Since $G$ is smooth, the $(A/(f^{1/n_1}_1,  f_2, \dotsc, f_m))$-module 
\[
\Hom_{A/(f_1,\, \dotsc,\, f_m )} (e^*(\Omega^1_{G/(A/(f_1,\, \dotsc,\, f_m))}), J)
\]
is finite projective. Consequently, by decreasing induction on $t$ and the regularity of the sequence $a_1, \dotsc, a_d$, for $i < d - t$ we have
\[
\tst H^i_I\p{A/(f_1,\, \dotsc,\, f_m ), \p{\Hom (e^*(\Omega^1_{G/(A/(f_1,\, \dotsc,\, f_m))}), J)}\!/(a_1, \dotsc, a_t)} \cong 0.
\]
The $t = 0$ case of this vanishing and the preceding short exact sequence then imply that
\[
H^i_I(A/(f_1,\, \dotsc,\, f_m ), G) \ra H^i_I(A/(f^{(n_1 - 1)/n_1}_1,\, f_2,\, \dotsc,\, f_m ), G) 
\]
is injective for $i < d$ and bijective for $i < d - 1$. By repeating this argument for 
\[
\Spec(A/(f^{(n_1 - 2)/n_1}_1, f_2,\, \dotsc,\, f_m)) \hra \Spec(A/(f^{(n_1 - 1)/n_1}_1, f_2,\, \dotsc,\, f_m ))
\]
and so on, we eventually replace $f_1$ by $f_1^{1/n_1}$ and then conclude by induction on $m$.
\epf



Another useful reduction is the following passage to a cover (compare with \cite{brauer-purity}*{Remark 2.7}). %We state in a setting adapted to our needs, although it is evidently a special case of a much more general statement. 

\blem \lab{finite-cover} 
Let $(R, \fm) \ra (R', \fm')$ be a finite, flat, local map of Noetherian local rings that are %with $R$ and $R'/\fm R'$ 
complete intersections, let $n < \dim(R)$ be an integer, and let $G$ be a commutative, finite, flat $R$-group. If for each local ring $(A, \fq)$ at some maximal ideal of some finite nonempty self-product $R' \tensor_R \dotsc \tensor_R R'$ Theorem~\uref{main} holds for $A$ and $G$ in cohomological degrees $< n$ in the sense that
\[
H^i_\fq(A, G) \cong 0 \qxq{for} i < n, 
\]
then 
\[
H^i_\fm(R, G) \cong 0 \qxq{for} i < n \qxq{and} H^n_\fm(R, G) \hra H^n_{\fm'}(R', G).
\]
%and $i \in \bZ$ such that
%\[%be \lab{CU-ass}
%H^{i'}_\fm(R' \tensor_R \dotsc \tensor_R R', G) = 0 \q \x{for all} \q i' < i   \qxq{and all finite } ,
%\]%ee
%with respect to the fppf cover $R'/R$, 
%we have
\elem

\bpf
The maximal ideals $\fq$ in question are the primes above $\fm$. Each $A$ is of dimension $\dim(R)$ and, by \eqref{vdim-ineq}--\eqref{vdim-add} (or by \cite{SP}*{Lemma~\href{https://stacks.math.columbia.edu/tag/09Q7}{09Q7}%, \href{https://stacks.math.columbia.edu/tag/069I}{069I}
}), is a complete intersection. Thus, the assumption gives 
\[
H^i_\fm(R' \tensor_R \dotsc \tensor_R R', G) \cong 0 \qxq{for} i < n.
\]
It then suffices to use the spectral sequence 
\[
E_2^{ij} = H^i(R'/R, H^j_\fm(-, G)) \Ra H^{i + j}_\fm(R, G)
\]
that results from fppf descent for $R'\mapsto R\Gamma_\fm(R',G)$ (and that could also be derived by choosing an injective resolution of $G$ and considering the \v{C}ech complexes of its terms with respect to $R'/R$).
\epf

%\beg \lab{finite-cover-rem} 
%The assumption that $R$ and $R'/\fm R'$ be complete intersections holds when $(R, \fm) \ra (R', \fm')$ is the reduction of a finite, flat, local map $(\wt{R}, \wt{\fm}) \ra (\wt{R}', \wt{\fm}')$ of regular local rings by an $\wt{R}$-regular sequence $f_1, \dotsc, f_n \in \wt{\fm}$ (see \cite{SP}*{\href{https://stacks.math.columbia.edu/tag/0E9K}{0E9K}}).
%implies that each self-product $R' \tensor_R \dotsc \tensor_R R'$ is also a ). Moreover, it 
%\eeg




\bprop \label{thm:main-pos-char}
\source{purity-flat-cohomology}
For a complete, Noetherian, local $\bF_p$-algebra $(R, \fm)$ that is a complete intersection with a perfect residue field $k$ and a commutative, finite, flat $R$-group $G$ of $p$-power order,
\[
H^i_\fm(R, G) \cong 0 \qxq{for} i < \dim(R).
\]
\eprop


\bpf
%Setting $k \ce R/\fm$, 
We use induction on $i$ simultaneously for all $R$. By the Cohen structure theorem (see \S\ref{conv}), 
\[
R \simeq k\llb t_1, \dotsc, t_N \rrb/(f_1, \dotsc, f_m)
\]
for a $k\llb t_1, \dotsc, t_N \rrb$-regular sequence $f_1, \dotsc, f_m$. By \Cref{finite-cover}, %(with \Cref{finite-cover-rem}) 
we may pass to the limit of the rings $k\llb t_1^{1/p^j}, \dotsc, t_N^{1/p^j} \rrb/(f_1, \dotsc, f_m)$ %for $N > 0$ 
to reduce to
\[
H^i_{(t_1,\, \dotsc,\, t_N)}(k\llb t_1^{1/p^{\infty}}, \dotsc, t_N^{1/p^{\infty}} \rrb/(f_1, \dotsc, f_m), G) \overset{?}{\cong} 0 \qxq{for} i < N - m.
\]
The ring $k\llb t_1^{1/p^\infty}, \dotsc, t_N^{1/p^\infty} \rrb$ is perfect, so %contains compatible systems of $p$-power roots $f_i^{1/p^\infty}$. Thus, 
\Cref{mod-out-by-more} and a limit argument reduce to 
\[%be \label{eqn:perfect-want-vanish}
\source{purity-flat-cohomology}
H^i_{(t_1,\, \dotsc,\, t_N)}(k\llb t_1^{1/p^{\infty}}, \dotsc, t_N^{1/p^{\infty}} \rrb/(f_1^{1/p^{\infty}}, \dotsc, f_m^{1/p^{\infty}}), G) \overset{?}{\cong} 0 \qxq{for} i < N - m.
\]%ee
Since 
\[
R_\infty \ce k\llb t_1^{1/p^{\infty}}, \dotsc, t_N^{1/p^{\infty}} \rrb/(f_1^{1/p^{\infty}}, \dotsc, f_m^{1/p^{\infty}})
\]
is a perfect $\bF_p$-algebra, \eqref{crystalline-cor} reduces us to
\[
H^i_{(p,\, t_1,\, \dotsc,\, t_N)}(W(R_\infty), \bM(G)) \overset{?}{\cong} 0 \qxq{for} i < N - m. 
\] 
Since the $W(R_\infty)$-module $\bM(G)$ is of projective dimension $\le 1$ (see \S\ref{Gab-eq}), this, in turn, reduces to 
\[
H^i_{(p,\, t_1,\, \dotsc,\, t_N)}(W(R_\infty), W(R_\infty)) \overset{?}{\cong} 0 \qxq{for} i < N - m + 1. %\qxq{and a free $W(R_\infty)$-module} M. 
\] 
By \cite{SP}*{Lemma~\href{https://stacks.math.columbia.edu/tag/07DV}{07DV}}, an $R$-regular sequence $a_1, \dotsc, a_{N - m} \in \fm$ is $R_\infty$-regular, so the sequence $a_0 \ce p$, $a_1$, $\dotsc$, $a_{\dim(R)}$ is $W(R_\infty)$-regular. Decreasing induction on $j$ then gives the sufficient
\[
H^i_{(p,\, t_1,\, \dotsc,\, t_N)}(W(R_\infty), W(R_\infty)/(a_0, \dotsc, a_j)) \cong 0 
\]
for $i < N - m - j$ and $-1 \le j \le N - m$.
\epf








%\temp{
%\blem \label{fppf-syn-comp}
%For a scheme $X$ and a commutative $X$-group scheme $G$ that is either smooth or finite locally free,
%\be \label{FSC-eq}
%R\Gamma_\syn(X, G) \isomto R\Gamma_\fppf(X, G).
%\ee
%\elem

%\bpf
%We let $\eps \colon X_\fppf \ra X_\syn$ and $\eps' \colon X_\syn \ra X_\et$ be the indicated morphisms of sites. %, we need to argue that $R^i \eps_*(G) = 0$ for $i > 0$. 
%For smooth $G$, by \cite{Gro68c}*{11.1, 11.2, and 11.6}, which were written for the fppf site but work the same way for the syntomic one,\footnote{\rv{In fact, the criterion \cite{Gro68c}*{11.2 (ii)} for the fppf site contains its syntomic version as a special case. }} we have $R^i(\eps' \circ \eps)_*(G) = 0$ and $R^i\eps'_*(G) = 0$ for $i > 0$, and consequently also
%\[
%R\Gamma_\et(X, G) \isomto R\Gamma_\syn(X, G) \qxq{and} R\Gamma_\et(X, G) \isomto R\Gamma_\fppf(X, G).
%\]
%This settles the case of a smooth $G$. For finite locally free $G$, by using the B\'{e}gueri resolution \eqref{Begueri-0} we obtain \eqref{FSC-eq} in cohomological degrees $i > 1$. For $i = 1$, both $H^1_\syn(X, G)$ and $H^1_\fppf(X, G)$ classify $G$-torsors because the latter are syntomic over $X$ (by \cite{SGA3Inew}*{III, 4.15} applied to the closed immersion $G \hra \Res_{G^*/X}(\bG_m)$ of Cartier duality, the structure map $G \ra X$ is syntomic and, by \cite{SP}*{\href{https://stacks.math.columbia.edu/tag/02VK}{02VK}}, every $G$-torsor inherits this property). In degree $i = 0$ the claim is evident.
%\epf
%}



















%%%%%%%%%%%%%%%%%%%%%%%%%%%%%%%%

\csub[Finite, locally free groups of $p$-power order over perfectoids]\label{section:Dieudonne-mixed-characteristic}
\source{purity-flat-cohomology}



The classification of commutative, finite, locally free groups of $p$-power order over perfect $\bF_p$-algebras %reviewed in \S\ref{section:Dieudonne-positive-characteristic} 
was extended to perfectoid $\bZ_p$-algebras by Lau \cite{Lau18} in the case $p > 2$ and by the second-named author \cite{SW19}*{Appendix to Lecture~17} in general by using ideas from integral $p$-adic Hodge theory. In \cite{ALB20}, Ansch\"utz--Le Bras drew a parallel to the crystalline theory %in  positive characteristic 
by relating this classification to the prismatic point of view. We will use these results for formulating (and proving) the general case of the key formula \eqref{eqn:key-formula}, so we now review them and include some relevant for us complements.

\bpp[Prismatic Dieudonn\'{e} modules over perfectoid rings] \label{pp:prismatic-review}\label{thm:dieudonne}
\source{purity-flat-cohomology}
For a perfectoid $\bZ_p$-algebra $A$ and a fixed generator $\xi$ of $\Ker(\bA_\Inf(A) \surjects A)$, by \cite{ALB20}*{Theorem 5.4} (which builds on \cite{SW19}*{Theorem~17.5.2}), there is a covariant, compatible with base change % This aspect is not stated super clearly in ALB19, but it is part of SW19, 17.5.2.
equivalence of categories
\[
G \mapsto \bM(G) \ce \Ext^1_{A_\Prism}(G^*, \mathcal O_\Prism)
\]
from the category of commutative, finite, locally free $A$-groups $G$ of $p$-power order to the category of finitely presented, $p$-power torsion $\bA_\Inf(A)$-modules $\bM$ of projective dimension $\le 1$ equipped with Frobenius (resp.,~inverse Frobenius) semilinear maps 
\[
F \colon \bM \ra \bM \ \ \x{(resp.,}\ \ V \colon \bM \ra \bM) \ \ \x{that satisfy} \ \ FV = \Frob(\xi) \ \ \x{and}\ \ VF = \xi.
\]
The $\Ext^1$ above is in the absolute prismatic site of $A$ %(viewed as the oriented perfect prism $(\bA_\Inf(A), \xi)$) 
and, by \cite{BS19}*{Lemma~4.8}, is an $\bA_\Inf(A)$-module. The Frobenius $F$ is induced by the Frobenius of the prismatic structure sheaf $\cO_{\Prism}$, and the only role of $\xi$ is to define the Verschiebung $V$. By \cite{ALB20}*{Theorem 4.44 and the proof of Theorem~5.4}, in the case when $A$ is a perfect $\bF_p$-algebra, %, granted that we chose $\xi \ce p$, 
the functor $G \mapsto \bM(G)$ may be identified with its crystalline counterpart discussed in \S\ref{Gab-eq}. % There is the point of the choice of $\xi$ in positive characteristic, but two different choices give an autoidentification of the functor $G \mapsto \bM(G)$ with itself, and we bake this into the compatibility when we say `may be identified.' 
By construction, in the case when $G$ is the $p^n$-torsion of a $p$-divisible group, $\bM(G)$ is a finite projective $(\bA_\Inf(A)/p^n\bA_\Inf(A))$-module.
\epp


\beg \label{explicit-M(G)}
\source{purity-flat-cohomology}
By \cite{ALB20}*{Lemma 4.75}, the $\bA_\Inf(A)$-module that underlies $\bM(\mu_{p^n})$ is % we start from this case because the citation is more direct in this case. Also, note that the cited reference is about $\bZ/p^n\bZ$ but we are precomposing with Cartier duality.  
\[
\bA_\Inf(A)/p^n\bA_\Inf(A) \qxq{with} F = \Frob(-),\q V = \xi\cdot\Frob\i(-).
\]
Likewise, %by also using \cite{ALB19}*{5.1.7}, 
the $\bA_\Inf(A)$-module that underlies ${\mathbf{1}}_n \ce \bM(\bZ/p^n\bZ)$ is  
\[
\bA_\Inf(A)/p^n\bA_\Inf(A) \qxq{with} F = \Frob(\xi)\cdot\Frob(-), \q V = \Frob\i(-).
\] 
\eeg




\blem \label{MG-exact-2}
\source{purity-flat-cohomology}
In \uS\uref{pp:prismatic-review}, the functor $G \mapsto \bM(G)$ and its inverse preserve short exact sequences. 
\elem

\bpf
Cartier duality is exact, so the claim is part of \cite{ALB20}*{Theorem~5.4}. %(Alternatively, one could adapt the argument of \Cref{MG-exact} above.)
\epf


%\kestutis{Reduce, to the characteristic $p$ case by considering $A_\flat$. The exactness of the inverse can hopefully be done the same was as in the proof of \cite{Ber80}*{3.4.2}.}



Similarly to \Cref{lem:syntomic-low-coh-id}, we obtain the following description of low degree cohomology of $G$.
%
%, the low degree cohomology $H^i(A, G)$ is described as follows.

\bprop \label{Hom-Ext-describe}
\source{purity-flat-cohomology}
For a perfectoid $\bZ_p$-algebra $A$ and a commutative, finite, locally free $A$-group $G$ of $p$-power order, we have  functorial in $A$ and $G$ identifications
\begin{equation} \label{MG-fixed-pts}
\source{purity-flat-cohomology}
\tst   G(A) \cong \bM(G)^{V = 1} \qxq{and} H^1(A, G) \cong \bM(G)/(V - 1)\bM(G). 
\end{equation}
\eprop

\bpf
For any $p^n$ that kills $G$, we have 
\[
H^1(A, G) \cong \Ext^1(\underline{\bZ/p^n\bZ}_A, G)
\]
(extensions of fppf $\bZ/p^n\bZ$-module sheaves), so the same argument as for \Cref{lem:syntomic-low-coh-id} (with \S\ref{pp:prismatic-review}, \Cref{explicit-M(G)}, and \Cref{MG-exact-2} in place of \S\ref{Gab-eq}, \Cref{Gab-eq-ex}, and \Cref{MG-exact}) gives the claim.
%For a $p^n$ that kills $G$, by , we have the functorial %in $A$ and $G$  
%identifications
%\[\ba
%G(A) &\cong \Hom(\underline{\bZ/p^n\bZ}_A, G) \cong \Hom_{\bA_\Inf(A),\, F,\, V}(\bb1_n, \bM(G)), \\
%H^1(A, G) &\cong \Ext^1(\underline{\bZ/p^n\bZ}_A, G) \cong \Ext^1_{\bA_\Inf(A),\, F,\, V}(\bb1_n, \bM(G)),
%\ea\]
%where the subscript ``$\bA_\Inf(A),\, F,\, V$'' indicates the target category of the functor $G \mapsto \bM(G)$ described in \S\ref{pp:prismatic-review} and the sheaves (resp.,~modules) representing the extensions in question are required to be killed by $p^n$. Moreover, by \Cref{explicit-M(G)}, giving a morphism $\bb1_n \ra \bM(G)$ amounts to specifying the image of $1 \in \bb1_n$ in $\bM(G)$---by $V$-equivariance, this image must be an $m \in \bM(G)^{V = 1}$, and any such $m$ works because $F$ is the multiplication by $\Frob(\xi)$ on $\bM(G)^{V = 1}$. %Thus, $G(A) \cong \bM(G)^{V = 1}$.
%For the remaining identification 
%\[
%\Ext^1_{\bA_\Inf(A),\, F,\, V}(\bb1_n, \bM(G)) \cong \bM(G)/(V - 1)\bM(G),
%\]
%we proceed analogously to the proof of \Cref{lem:syntomic-low-coh-id} (in which $\xi = p$): any $m \in \bM(G)$ gives an extension structure on $\bM(G) \oplus \bb1_n$ by requiring that the Verschiebung send $(0, 1)$ to $(m, 1)$ (so that the Frobenius then sends $(0, 1)$ to $(-F(m), \Frob(\xi))$), two such extensions are isomorphic if and only if the difference of the corresponding elements $m$ lies in $\bM(G)^{V = 1}$, and any extension is of this form for some $m$ because it splits as an extension of $(\bA_\Inf(A)/p^n\bA_\Inf(A))$-modules.
\epf






\brem
In \Cref{Hom-Ext-describe}, the $H^{\ge 2}(A, G)$ vanish: we will deduce this in \Cref{no-higher-coho} from its positive characteristic case of \Cref{cor:perfect-vanish} via the $p$-adic continuity formula of \S\ref{section:continuity-formula}.
\erem




We turn to analyzing the prismatic side of the key formula \eqref{eqn:key-formula}: we show that it satisfies $p$-complete arc descent in \Cref{M(G)-coho-descent} and then arc locally relate it to flat cohomology in \Cref{M(G)-local-analysis}. A key input to our %$p$-complete 
arc descent results is the following lemma that in essence restates \cite{BS19}*{Proposition~8.10}.

\blem \label{arc-sheaves}
\source{purity-flat-cohomology}
The following functors satisfy hyperdescent for those $p$-complete arc hypercovers whose terms are perfectoid $\bZ_p$-algebras\ucolon
\begin{equation*} \label{}
\source{purity-flat-cohomology}
A \mapsto \bA_\Inf(A), \q A \mapsto A, \qxq{and} A \mapsto A^\flat.%, \qxq{and} A \mapsto (A/p A)^\red.
\end{equation*}
\elem

\bpf
Since $\Ker(\theta\colon \bA_{\Inf}(A) \surjects %\xra{[x] \mapsto x^\sharp} 
A)$ and $\Ker(\bA_\Inf(A) \surjects A^\flat)$ are principal, generated by nonzerodivisors $\xi$ and $p$, and $\bA_\Inf(A)$ is $p$-adically complete, it suffices to treat $A \mapsto A^\flat$. Thus, fixing $A$ and letting $\wh{(-)}$ denote derived $p^\flat$-adic completion, \Cref{arc-tilt} (with \Cref{recognize}~\ref{R-a}) % The parenthetical reference is supposed to help justify that a levelwise tilt of a $p$-complete hypercover is a $p^\flat$-complete hypercover. We do not  reduce to descent at the outset because technically the functors are not defined on the objects of the Cech nerve of a cover (they only become defined after p-adic completions).
reduces us to showing that the functor $S \mapsto \wh{S}$ satisfies $p^\flat$-complete arc hyperdescent on the category of perfect $A^\flat$-algebras (by \S\ref{perfectoid-def} and \eqref{no-big-tor}, the tilts of perfectoid $A$-algebras are derived $p^\flat$-adically complete, so on them this functor agrees with $S \mapsto S$). Since the functor is bounded below (even concentrated in degree $0$), hyperdescent for it is equivalent to descent.  Moreover, it suffices to show arc descent: indeed, if $S \ra S'$ is a $p^\flat$-complete arc cover, then 
\[
\tst S \ra S' \times S[\frac 1{p^\flat}]
\]
is an arc cover and the functor %$S \mapsto \wh{S}$ 
has identical values on the two \v{C}ech nerves. We then instead consider the functor $S\mapsto \wh{S_\perf}$ defined on \emph{all} $A^\flat$-algebras, where 
\[
\tst S_\perf \ce \varinjlim_{s \mapsto s^p} S,
\]
and then reduce further to showing arc descent for $S \mapsto S_\perf$ on the category of all $A^\flat$-algebras $S$. By \cite{BS17}*{Theorem~4.1~(i), Proposition~4.5}, % certainly, $v$-cohomology satisfies $v$-descent, which is how we apply 4.1 (i) there.
the functor $S \mapsto S_\perf$ satisfies $v$-descent, so, by  \cite{BM20}*{Proposition~4.8 and its proof}, % The ``and its proof'' part is supposed to take care of the fact that there the statement is only made for functor defined on all qcqs schemes rather than just affine schemes.
it also satisfies arc descent. 
%For this, we follow the argument of \cite{BS19}*{8.9}. Namely, by , if $A \ra A'$ is a $p$-complete arc cover of perfectoids, then $A^\flat \ra A'^\flat \times A^\flat[\f 1{p^\flat}]$ is an arc cover of perfect $\bF_p$-algebras. Thus, it suffices to show that on perfect $\bF_p$-algebras the functor $S \mapsto S$ satisfies arc descent: indeed,
%By \cite{BM20}*{Thm.~1.6} with \cite{BS17}*{4.1 and 6.3}, 
%. Thus, on the category of perfect $A^\flat$-algebras the functor $B \mapsto \wh{B}$ satisfies arc descent, where the completion is $p^\flat$-adic and derived (equivalently, nonderived because the $(p^\flat)^\infty$-torsion of $B$ is bounded). Completion eliminates $A^\flat[\f 1{p^\flat}]$, and we obtain $p$-complete arc descent for $A \mapsto A^\flat$. 
\epf



\bprop \label{M(G)-coho-descent}
\source{purity-flat-cohomology}
Let $A$ be a perfectoid $\bZ_p$-algebra,  let   $Z \subset \Spec(A/pA)$ be a closed subset, let $a \in \bA_\Inf(A)$, and %that is $\varpi$-adically complete for a $\varpi \in A$ with $\varpi^p \mid p$, 
let $G$ be a commutative, finite, locally free $A$-group of $p$-power order. %on %the category of %$\varpi$-adically complete 
 The functors
\[
\tst A' \mapsto \bM(G_{A'})[\f 1a] \qxq{and} A' \mapsto R\Gamma_Z(\bA_\Inf(A'), \bM(G_{A'})) %\qx{satisfy  hyperdescent.}
\]
satisfy hyperdescent for those $p$-complete arc hypercovers whose terms are perfectoid $A$-algebras.
\eprop

\bpf
For any cosimplicial abelian group $M^\bullet$, the associated complex 
\[
M^0 \ra M^1 \ra \dotsc
\]
represents $R\lim_\Delta(M^\bullet)$, % See MO # 338498
so, since localization is exact, we may assume that $a = 1$. % This is slight
 %$\bM(G_{A'})[\f 1a] \cong \varinjlim\p{\bM(G_{A'}) \xra{a} \bM(G_{A'}) \xra{a} \dotsc}$, we may assume that $a = 1$. % 
By  \S\ref{pp:prismatic-review}, we have 
\[
\bM(G_{A'}) \cong \bM(G) \tensor_{\bA_\Inf(A)} \bA_\Inf(A')
\]
and there are finite projective $\bA_\Inf(A)$-modules $M_i$ that fit into an exact sequence 
\[
0 \ra M_0 \ra M_1 \ra \bM(G) \ra 0.
\]
Since $\bM(G)$ is $p$-power torsion and $\bA_\Inf(A')$ is $p$-torsion free, this sequence stays exact after base change to $\bA_\Inf(A')$. Thus, for the claim about $A' \mapsto \bM(G_{A'})$ it suffices %to resolve $\bM(G)$ by finite projective $\bA_\Inf(A)$-modules $M_i$ and 
note that, by \Cref{arc-sheaves}, the functors 
\[
A' \mapsto M_i \tensor_{\bA_\Inf(A)} \bA_\Inf(A')
\]
satisfy hyperdescent for those $p$-complete arc hypercovers whose terms are perfectoid $A$-algebras. For the claim about the functor
\[
A' \mapsto R\Gamma_Z(\bA_\Inf(A'), \bM(G_{A'})),
\]
we use the functorial triangle 
\be
\ba
\tst R\Gamma_Z(\bA_\Inf(A'), \bM(G_{A'})) \ra &R\Gamma(\bA_\Inf(A'), \bM(G_{A'})) \\ &\tst\ra R\varprojlim_z \p{R\Gamma(\bA_\Inf(A')[\f{1}{z}], \bM(G_{A'})[\f 1z])}
\ea
\ee
where $z$ ranges over the elements of $\bA_\Inf(A)$ that vanish on $Z$. This reduces us to the settled claim about $A' \mapsto \bM(G_{A'})[\f 1a]$ because 
\[
\tst R\Gamma(\bA_\Inf(A')[\f{1}{z}], \bM(G_{A'})[\f 1z]) \cong \bM(G_{A'})[\f 1z]. \qedhere
\]
\epf






The promised arc-local analysis uses the following 
%We conclude the section with \Cref{M(G)-local-analysis} that describes the arc local (or even $v$-local) properties of the prismatic Dieudonn\'{e} modules $\bM(G)$. These properties will be used in the proof of \Cref{thm:perfectoidsupport}; for their argument, we review a 
general lemma about modules on infinite products. 


\blem \label{product-lemma}
\source{purity-flat-cohomology}
For rings $\{R_i\}_{i \in I}$ and a finitely presented $(\prod_{i \in I} R_i)$-module $M$, we have %the module isomorphism
\begin{equation*}
\tst M \isomto \prod_{i \in I} (M \tensor_R R_i).
\end{equation*}
\elem

\bpf
Set $R\ce \prod_{i\in I} R_i$ and choose a resolution $R^n\to R^m\to M\to 0$. Both $-\otimes_R R_i$ and 
%is right exact (even exact -- $R_i$ is flat over $R$ as it splits off as a direct factor), and 
infinite products are exact, so the claim for $M$ reduces to the evident case of a finite %$M$ is finite 
free $R$-module. %modules $R^n$ and $R^m$, where it is clear.
\epf





%The preceding lemma immediately gives the following further property of the functor $G \mapsto \bM(G)$. 

\bprop \label{M(G)-products}
\source{purity-flat-cohomology}
Let $\{A_i \}_{i \in I}$ be perfectoid $\bZ_p$-algebras, let $A \ce \prod_{i \in I} A_i$ \up{which is perfectoid by Proposition \uref{recognize}~\uref{R-c}}, and let $G$ be a commutative, finite, locally free $A$-group of $p$-power order. We have
\[
\tst \bM(G) \isomto \prod_{i \in I} \bM(G_{A_i})
\]
over 
\[
\tst \bA_\Inf(A) \isomto \prod_{i \in I} \bA_\Inf(A_i)
\]
compatibly with $F$ and $V$, granted that we choose compatible $\xi$ in $\bA_\Inf(A)$ and $\bA_\Inf(A_i)$ to define $V$ \up{see \uS\uref{pp:prismatic-review}}.
\eprop

\bpf
Since $G \mapsto \bM(G)$ commutes with base change (see \S\ref{pp:prismatic-review}), compatibility with $F$ and $V$ is automatic and \Cref{product-lemma} gives the claim.
\epf






\bprop \label{M(G)-local-analysis}
\source{purity-flat-cohomology}
Let $\{A_i\}_{i \in I}$ be $p$-adically complete valuation rings  %of rank $\le 1$ 
of rank $\le 1$ such that  the fraction field $K_i$ of $A_i$ is algebraically closed, let $A \ce \prod_{i \in I} A_i$, let $K \ce \prod_{i \in I} K_i$, let $a \in A$, and let $G$ be a commutative, finite, locally free $A$-group of $p$-power order. The map
\[
\tst V - 1 \colon \bM(G)[\f 1{[a^\flat]}] \ra \bM(G)[\f 1{[a^\flat]}]
\]
is surjective \up{for any $\xi \in \bA_\Inf(A)$ used to define $V$ of $\bM(G)$, see \uS\uref{pp:prismatic-review}}, 
\begin{equation*}
\tst H^j(A[\f 1a], G) \cong 0 \qxq{for} j \ge 1,
\end{equation*}
%\begin{equation*}
%\tst \bM(G)^{V = 1} \surjects (\bM(G)[\f 1{a^\flat}])^{V = 1} \qxq{and}  \qx{are surjective}
%\end{equation*}
and there is a unique commutative~square
\[%begin{equation} \ba\label{M(G)-square}
\source{purity-flat-cohomology}
\tst \xymatrix@C=40pt{
G(A) \ar@{->>}[d]\ar[r]_-{\eqref{MG-fixed-pts}}^-{\sim} & \bM(G)^{V = 1} \ar@{->>}[d] \\
G(A[\f 1{a}]) \ar[r]^-{\sim} & (\bM(G)[\f 1{[a^\flat]}])^{V = 1},
}
\]%ea \end{equation}
in which the vertical arrows are bijective whenever $a$ is a nonzerodivisor. 
\eprop

\emph{Proof.}
%By \Cref{recognize}~\ref{R-c}, the ring $A$ is perfectoid. 
In the statement, $A$ is perfectoid by \eqref{p-oid-mod-pi} and \Cref{recognize}~\ref{R-c}, the element $a^\flat \in A^\flat \cong \prod_{i \in I} A_i^\flat$ is a system of compatible $p$-power roots of $a \in A$ (see \eqref{monoid-id}), and the map  
\[
\tst V \colon \bM(G)[\f 1{[a^\flat]}] \ra \bM(G)[\f 1{[a^\flat]}]
\]
is defined by 
\[
\tst V(\f{m}{[(a^\flat)^n]}) \ce \f{V(m)}{[(a^\flat)^{n/p}]}.
\]
By decomposing $A$ into subproducts, we may replace $I$ by parts of a finite partition of $I$. Thus, since the case $a = 0$ is clear, we assume that $a$ is a nonzerodivisor. In this case, since, by \Cref{tilt-summary}, each $A_i^\flat$ is a valuation ring, $a^\flat$ is also a nonzerodivisor, and the uniqueness aspect %of \eqref{M(G)-square} 
will follow from the bijectivity of the vertical maps. Moreover, $A \hra A[\f 1a] \hra K$, so, since $G(A) \isomto G(K)$ by the valuative criterion of properness, %\cite{EGAI}*{9.5.6} bootstraps the desired bijectivity 
also $G(A) \isomto G(A[\f 1a])$.
Due to the local structure of $A[\f 1a]$ described in \Cref{lem:ultrafilters}, the valuative of properness applied locally on $\Spec(A[\f{1}{a}])$ also ensures that $X(A[\f 1a]) \isomto X(K)$
for every $G_{A[\f{1}{a}]}$-torsor $X$. Since $K$ is a product of algebraically closed fields, %and $X$ is affine, 
this shows that %every $G_{A[\f{1}{a}]}$-torsor is trivial, in other words, that 
\[
\tst H^1(A[\f 1a], G) \cong 0.
\]
In contrast, \Cref{lem:ultrafilters} and the B\'{e}gueri sequence \eqref{Begueri-0} show that 
\[
\tst H^j(A[\f 1a], G) \cong 0 \qxq{for} j \ge 2.
\]
The remaining claims concern $\bM(G)$, and we use \cite{BBM82}*{th\'{e}or\`{e}me~3.1.1} to cover $\Spec(A)$ by finitely many affine opens $\Spec(A')$ over which $G$ embeds into a $p$-divisible group. Since each $A_i$ is local, the map $A \ra A_i$ factors through some $A \ra A'$, so we subdivide $I$ to assume that there is an exact sequence 
\begin{equation*} \label{}
\source{purity-flat-cohomology}
0 \ra G \ra G' \ra G'' \ra 0
\end{equation*}
in which $G'$ is a truncated $p$-divisible group. Since 
\[
G(A) \isomto \bM(G)^{V = 1} \qxq{and} H^1(A, G) \cong 0,
\]
%Due to the bijectivity of the top horizontal map in the square %\eqref{M(G)-square}  and the vanishing 
the functor 
\[
\tst G \mapsto \bM(G)^{V = 1} \qx{is exact}
\]
and, by \Cref{MG-exact-2}, the functor 
\[
\tst G \mapsto (\bM(G)[\f 1{[a^\flat]}])^{V = 1} \qx{is left exact.}
\]
 Thus, by snake lemma, for the injectivity and then also the surjectivity of the right vertical map in the diagram, %\eqref{M(G)-square}, 
we may replace $G$ by $G'$ to assume that $G$ is a truncated $p$-divisible group. By then $p$-adically filtering $G$ and again using the snake lemma, we may assume that $G$ is also killed by $p$. In this case, $\bM(G)$ is a projective $A^\flat$-module (see \S\ref{pp:prismatic-review}) and the right vertical map is injective because $a^\flat$ is a nonzerodivisor in $A^\flat$. %\overset{\x{\ref{recognize}~\ref{R-c}}}{\cong} \prod_{i\in I} A_i^\flat$ (each $A_i^\flat$ is a valuation ring of rank $\le 1$, see \Cref{tilt-summary}). 
For its surjectivity, fix an 
\[
\tst m \in (\bM(G)[\f 1{a^\flat}])^{V = 1}.
\]
By the $\Frob\i$-semilinearity of $V$, if $bm \in \bM(G)$ for a $b \in A^\flat$, then also $b^{1/p}m \in \bM(G)$. Such elements of $A^\flat[\f 1{a^\flat}]$ lie in $A^\flat$, so, since $\bM(G)$ is a direct summand of a finite free $A^\flat$-module, $m \in \bM(G)$. %, and hence $(\bM(G)[\f 1{a^\flat}])^{V = 1} \subset \bM(G)^{V = 1}$.
%on  to conclude that $(a^\flat)^{1/p^n}((\bM(G)[\f 1{a^\flat}])^{V = 1}) \subset \bM(G)$ for every $n > 0$. Since each $A_i^\flat$ is of rank $\le 1$, any element of $A^\flat[\f 1{a^\flat}]$ that is brought into $A^\flat$ by scaling by every $(a^\flat)^{1/p^n}$ must lie in $A^\flat$, so, by , we get 
%In fact, since $\prod_{i \in I} K_i^\flat$, where $K_i^\flat$ is the fraction field of $A_i^\flat$, is the filtered direct limit of the $A^\flat[\f 1{a^\flat}]$ as $a^\flat$ ranges over the nonzerodivisors in $A^\flat$, the argument shows that even $\bM(G)^{V -1} \isomto (\bM(G) \tensor_{A^\flat} (\prod_{i \in I} K_i^\flat))^{V = 1}$ for truncated $p$-divisible groups $G$ killed by $p$. 
  

For the remaining surjectivity of 
\[
\tst \tst V - 1 \colon \bM(G)[\f 1{[a^\flat]}] \surjects \bM(G)[\f 1{[a^\flat]}],
\]
by \Cref{MG-exact-2}, the bijectivity of the right vertical arrow, %\eqref{M(G)-square}, 
and the snake lemma, the functor 
\[
\tst G \mapsto \bM(G)[\f 1{[a^\flat]}]/(V - 1)(\bM(G)[\f 1{[a^\flat]}]) \qx{is exact.}
\]
Thus, as before, we may first assume that $G$ is a truncated $p$-divisible group and then that it is also killed by $p$, so that $\bM(G)$ is a finite projective $A^\flat$-module. We then first claim that $V - 1$ is surjective after passing to the limit over all the nonzerodivisors   $a^\flat$.

\bcl \label{first-claim}
\source{purity-flat-cohomology}
Letting $K_i^\flat$ be the fraction field of $A_i^\flat$, the map $V - 1$ is surjective on
\[
\tst \bM(G) \tensor_{A^\flat} (\prod_{i \in I} K_i^\flat) \overset{\ref{product-lemma}}{\cong}  \prod_{i \in I} (\bM(G_{A_i}) \tensor_{A_i^\flat} K_i^\flat).
\] 
\ecl

\bpf
%Namely, by , we have $\bM(G) \tensor_{A^\flat} (\prod_{i \in I} K_i^\flat) $, so for the claim 
We may assume that $I$ is a singleton $\{ i \}$. If $A_i$ is an $\bF_p$-algebra, then 
\[
\bM(G) \tensor_{A_i^\flat} K_i^\flat \cong \bM(G_{K_i})
\]
and $V - 1$ is surjective by \eqref{MG-fixed-pts}. %because $H^1(K_i, G) = 0$. 
Otherwise $\xi$ is a nonzerodivisor in $A^\flat$, so both $F$ and $V$ are bijective on 
\[
M \ce \bM(G_{A_i}) \tensor_{A_i^\flat} K_i^\flat.
\]
Thus, by \cite{Kat73}*{Proposition 4.1.1 and its proof}, $M^{F = 1}$ is a finite $\bF_p$-module and $M \cong M^{F = 1} \tensor_{\bF_p} K_i^\flat$ compatibly with the Frobenius. By choosing an $\bF_p$-basis for $M^{F = 1}$, the desired surjectivity of $V - 1$ on $M$ then amounts to the surjectivity of $\xi \cdot \Frob\i(*) - *$ on $K_i^\flat$, equivalently, of $\xi \cdot * - \Frob(*)$. The latter translates to the solvability in $K_i^\flat$ of any equation $x^p - \xi x + b = 0$  with $b \in K_i^\flat$, which follows from the algebraic closedness of $K_i^\flat$ (see \Cref{tilt-summary}).%, the field $K_i^\flat$ is algebraically closed, so such equations are indeed solvable in $K_i^\flat$.
\epf 


%Keeping the assumption that $G$ is a truncated $p$-divisible group killed by $p$, 
Thanks to \Cref{first-claim}, 
it remains to show that if an $m \in \bM(G)[\f 1{a^\flat}]$ is of the form $V(x) - x$ for some $x \in \bM(G) \tensor_{A^\flat} (\prod_{i \in I} K_i^\flat)$, then $x \in \bM(G)[\f 1{a^\flat}]$. For this, it suffices to show that $x$ lies in each stalk of $\bM(G)$ at a variable point of $\Spec(A[\f 1{a^\flat}])$. %, so we may work locally on $\Spec(A[\f 1{a^\flat}])$. 
By \Cref{lem:ultrafilters}, each local ring of $\Spec(A[\f 1{a^\flat}])$ is a valuation ring whose fraction field is a localization of $K$, so we are reduced to the following claim.

\bcl
For a perfect $\bF_p$-algebra $W$ that is a valuation ring with fraction field $L$, a finite free $W$-module $M$, and a $\Frob\i$-semilinear $V \colon M \ra M$, any $x \in M_L$ with $V(x) - x \in M$ lies in $M$.
\ecl

\emph{Proof. }
In the statement, the map $V$ on $M_L$ is defined by the same formula as in the beginning of the proof of \Cref{M(G)-local-analysis}. Suppose that $w \in W$ is such that $wx \in M$, so that $x = \f{m}{w}$ for some $m \in M$. Then, since $V(x) - x = \f{V(m)}{w^{1/p}} - x$ lies in $M$, we get that also $w^{1/p}x \in M$. However, since $M$ is finite free and $w$ is arbitrary subject to $wx \in M$, this means that $x \in M$.
%, any $x \in M_L$ such that $wx \in M$ implies $w^{1/p}x \in M$ must lie in $M$. 
\QEDD
















%\bprop \label{M(G)-invert-pi}
%For a perfectoid $A$ that is $\varpi$-adically complete for a $\varpi \in A$ with $\varpi^p \mid p$, the following functors on the category of commutative, finite, locally free, $p$-power order, $A[\f 1\varpi]$-\'{e}tale $A$-groups are naturally isomorphic functorially in $A$ and compatibly with replacing $\varpi$ by a multiple\ucolon
%\[
%\tst G \mapsto \bM(G)[\f 1{\varpi^\flat}] \qxq{and} G \mapsto \bM((G_{A[\f  1\varpi]})^\flat).
%\]
%\eprop

%\bpf

%\epf



%\bcor \label{G(A)-invert-pi}
%For a perfectoid $A$ that is $\varpi$-adically complete for a $\varpi \in A$ with $\varpi^p \mid p$ and a commutative, finite, locally free, $p$-power order, $A[\f 1\varpi]$-\'{e}tale $A$-group $G$,
%\begin{equation*}
%\tst G(A[\f 1\varpi]) \cong (\bM(G)[\f 1{\varpi^\flat}])^{V - 1} \qxq{compatibly with} G(A) \overset{\ref{MG-fixed-pts}}{\cong} \bM(G)^{V - 1}, 
%\end{equation*}
%functorially in $G$ and $A$, and compatibly with making $\varpi$ more divisible. 
%\ecor

%\bpf

%\epf







%The arc descent property of $\bM(G)$ of \Cref{M(G)-coho-descent} will reduce us to products of valuation rings. 














%%%%%%%%%%%%%%%%%%%%%%%%%%%%%%%%%%%%%%%%%%%%%%
